\section{Conclusion}
Language models have developed increasingly good ``procedural'' knowledge of what should be discussed and how it should be presented, but often struggle to memorize ``factoid'' knowledge and produce unsubstantiated claims.
We proposed \approach{}, a framework for revising such claims to make them attributable to the researched evidence.
From experiments on text passages generated by different models on various domains, we showed that \approach{} can revise the passages to improve attribution while preserving other desirable properties such as writing style or structure.
Furthermore, \approach{} sits on top of existing generation models without needing to re-design or re-train LMs.

Major headroom still remains, as discussed in Section~\ref{sec:analysis} and the Limitations section. We hope our analysis of \approach{} would help with developing new approaches for integrating attribution to LMs.


\section{Limitations}

\paragraph{Limitations of our task definition}
Depending on the application, attribution and preservation may not deserve equal weight. For instance, if there are multiple acceptable options for the output, such as in a dialog system, we might trade-off preservation for attribution, similar to how LaMDA behaves in our experiments.

Our evaluation metrics also do not measure all aspects of attribution. For instance, some sentences are self-evident and do not require attribution (e.g., \emph{``I agree.''}) but would be penalized in our evaluation.
It is also necessary to note that linguistic assertions have varying scope: for example, there is a difference between \emph{``Frozen is a scary movie''} and \emph{``I got scared watching Frozen''} --- while expressing a similar sentiment, the former makes a more general statement that many would disagree with, while the latter is scoped to the speaker's own experience. In some applications, one could even argue that the latter case does not require attribution, since the speaker is their own source-of-truth. In addition to varying scope, utterances can also make assertions with varying levels of directness. For example, according to standard linguistics, \emph{``John ate some of the cookies''} yields the implicature that John did not eat \emph{all} of the cookies, even though it is not logically entailed. This raises the question of which implicatures or implied assertions should be detected and attributed, which should be explored in future work. For more nuances, we refer to \citet{Rashkin2021AIS}.

For preservation, we wish to explore other properties that should be preserved, such as discourse or logical coherence. Additionally, if the input text passage is completely misguided or flawed, it can be difficult to revise the text without significant changes, which would be heavily penalized by the current metrics.

\paragraph{Limitations of our model}
While we aspire to improve attribution for arbitrary text, it is clear that \approach{} is not yet fully general. For example, the current implementation of \approach{} would not be well-prepared to edit poetry (where preserving rhyme matters) or long documents, primarily because we do not provide examples of such inputs in our few-shot LLM prompts. However, we do believe that future developers may be able to quickly adapt \approach{} to such tasks by simply changing the prompts. Second, \approach{} tends to preserve rather than delete claims that it cannot attribute. Some of these claims genuinely do not require attribution, but others are hallucination and should be removed. Judging whether a claim requires attribution can be subjective and challenging. Finally, our model is computationally costly, since it is based on prompting a large language model. One potential solution is to leverage recent synthetic data generation recipes to train a smaller model \cite{lee2021neural,Schick2022PEER}.

\section{Ethical considerations}

\paragraph{Partial attribution}
When \approach{} is not 100\% successful in making text consistent with retrieved evidence, the revised text will be partially attributed.
One could identify unattributed parts using either the automated attribution score ($\attrais{}$) or the relevance scores used to generate the attribution report (Section~\ref{sec:attribution-report}).
Such information should be presented to avoid misleading readers into thinking that the entire revision is attributed.

\paragraph{Evidence trustworthiness}
\approach{} seeks to improve attribution for the output of any generative model. However, even if \approach{} can attribute content to a particular source, the user must still consider whether the source itself is trustworthy. Even for sources that are traditionally considered ``authoritative'' (such as an encyclopedia), there may still be factual inaccuracies or biases. This work does not address the question of whether a source is trustworthy, or the related topic of misinformation.
While we do not provide a means for judging trustworthiness, the design of \approach{} does allow for the research stage to restrict its search over a user-specified corpus, based on what the user deems trustworthy.

\paragraph{Conflicting evidence}
There is also the possibility that some content may be simultaneously supported by certain sources, while contradicted by others. This can easily occur for content involving subjective or imprecise claims. The current implementation and evaluation for \approach{} does not explicitly address this issue --- we adopted a ``permissive'' definition of attribution, where we consider content to be attributed if there exists any source that supports it. For some applications, a more restrictive definition that requires both existence of supporting sources and absence of contradicting sources would be needed.

